\QHeading{第1题　口算。}
\par\noindent


\begin{oralgrid}
\OralItem{1}{26\times70}{1820}
\OralItem{2}{449+21}{470}
\OralItem{3}{201-96}{105}
\OralItem{4}{545-235}{310}
\OralItem{5}{340\div20}{17}
\OralItem{6}{785-237}{548}
\OralItem{7}{25\times40}{1000}
\OralItem{8}{820-570}{250}
\OralItem{9}{510\div17}{30}
\OralItem{10}{920\div4}{230}
\OralItem{11}{24\times20}{480}
\OralItem{12}{125\times20}{2500}\end{oralgrid}
\par

\QHeading{口算思路补充（对应上表题号）}

\Answer{上面列出得数，下面按小题给出不同计算思路。计数单位法适用于小数，整数题可比较凑整、拆分或商不变等方法。}
\MethodTitle{第1小题：$26\times 70$}

\MethodBlock{口算思路一：按数位拆分计算}{把一个因数拆成整十、整百数与余下部分之和。}{\begin{aligned}[t]&\phantom{={}}26\times 70\\
&=70\times (20+6)\\
&=70\times 20+70\times 6\\
&=1400+420\\
&=1820\end{aligned}}
\MethodTitle{第3小题：$201-96$}

\MethodBlock{口算思路一：补差法}{从减数先补到一个整数，再补到被减数，两段差相加。}{\begin{aligned}[t]&\phantom{={}}201-96\\
&=(100-96)+(201-100)\\
&=4+101\\
&=105\end{aligned}}
\MethodTitle{第4小题：$545-235$}

\MethodBlock{口算思路一：补差法}{从减数先补到一个整数，再补到被减数，两段差相加。}{\begin{aligned}[t]&\phantom{={}}545-235\\
&=(240-235)+(545-240)\\
&=5+305\\
&=310\end{aligned}}
\MethodTitle{第5小题：$340\div 20$}

\MethodBlock{口算思路一：同时去掉相同个数的0}{}{\begin{aligned}[t]&\phantom{={}}340\div 20\\
&=(340\div 10)\div (20\div 10)\\
&=34\div 2\\
&=17\end{aligned}}

\MethodBlock{口算思路二：拆分被除数}{}{\begin{aligned}[t]&\phantom{={}}340\div 20\\
&=(200+140)\div 20\\
&=200\div 20+140\div 20\\
&=10+7\\
&=17\end{aligned}}
\MethodTitle{第6小题：$785-237$}

\MethodBlock{口算思路一：补差法}{从减数先补到一个整数，再补到被减数，两段差相加。}{\begin{aligned}[t]&\phantom{={}}785-237\\
&=(240-237)+(785-240)\\
&=3+545\\
&=548\end{aligned}}
\MethodTitle{第7小题：$25\times 40$}

\MethodBlock{口算思路一：拆因数，凑整计算}{把一个因数拆成两个因数，先凑成整百或整千。}{\begin{aligned}[t]&\phantom{={}}25\times 40\\
&=25\times (4\times 10)\\
&=(25\times 4)\times 10\\
&=100\times 10\\
&=1000\end{aligned}}

\MethodBlock{口算思路二：按数位拆分计算}{把一个因数拆成整十、整百数与余下部分之和。}{\begin{aligned}[t]&\phantom{={}}25\times 40\\
&=40\times (20+5)\\
&=40\times 20+40\times 5\\
&=800+200\\
&=1000\end{aligned}}
\MethodTitle{第8小题：$820-570$}

\MethodBlock{口算思路一：补差法}{从减数先补到一个整数，再补到被减数，两段差相加。}{\begin{aligned}[t]&\phantom{={}}820-570\\
&=(580-570)+(820-580)\\
&=10+240\\
&=250\end{aligned}}
\MethodTitle{第10小题：$920\div 4$}

\MethodBlock{口算思路一：利用商不变的性质}{被除数和除数同时乘同一个非零数，商不变。}{\begin{aligned}[t]&\phantom{={}}920\div 4\\
&=(920\times 5)\div (4\times 5)\\
&=4600\div 20\\
&=230\end{aligned}}

\MethodBlock{口算思路二：拆分除数，连续相除}{}{\begin{aligned}[t]&\phantom{={}}920\div 4\\
&=920\div (2\times 2)\\
&=920\div 2\div 2\\
&=460\div 2\\
&=230\end{aligned}}
\MethodTitle{第11小题：$24\times 20$}

\MethodBlock{口算思路一：按数位拆分计算}{把一个因数拆成整十、整百数与余下部分之和。}{\begin{aligned}[t]&\phantom{={}}24\times 20\\
&=20\times (20+4)\\
&=20\times 20+20\times 4\\
&=400+80\\
&=480\end{aligned}}
\MethodTitle{第12小题：$125\times 20$}

\MethodBlock{口算思路一：按数位拆分计算}{把一个因数拆成整十、整百数与余下部分之和。}{\begin{aligned}[t]&\phantom{={}}125\times 20\\
&=20\times (100+25)\\
&=20\times 100+20\times 25\\
&=2000+500\\
&=2500\end{aligned}}


\par\noindent

\QHeading{第2题　脱式计算。}

\CalcBlock{1}{\begin{aligned}[t]&\phantom{={}}\color{black}25\times[(300-260)\div5]\\&\color{answerred}=25\times[40\div5]\\&\color{answerred}=25\times8\\&\color{answerred}=200\end{aligned}}
\par

\CalcBlock{2}{\begin{aligned}[t]&\phantom{={}}\color{black}156\div[480\div(276-196)]\\&\color{answerred}=156\div[480\div80]\\&\color{answerred}=156\div6\\&\color{answerred}=26\end{aligned}}
\CalcBlock{3}{\begin{aligned}[t]&\phantom{={}}\color{black}[288\div(295-289)]\times74\\&\color{answerred}=[288\div6]\times74\\&\color{answerred}=48\times74\\&\color{answerred}=3552\end{aligned}}
\CalcBlock{4}{\begin{aligned}[t]&\phantom{={}}\color{black}420\div(70-40)\times5\\&\color{answerred}=420\div30\times5\\&\color{answerred}=14\times5\\&\color{answerred}=70\end{aligned}}
\CalcBlock[方法一]{5}{\begin{aligned}[t]&\phantom{={}}\color{black}93\times[188-(78+53)]\\&\color{answerred}=93\times[188-131]\\&\color{answerred}=93\times57\\&\color{answerred}=5301\end{aligned}}
\MethodBlock{方法二：乘法分配律}{把括号中的两部分分别相乘，再相加或相减。}{\begin{aligned}[t]&\phantom{={}}93\times [188-(78+53)]\\
&=93\times 188-93\times (78+53)\\
&=93\times 188-93\times 131\\
&=17484-12183\\
&=5301\end{aligned}}

\CalcBlock[方法一]{6}{\begin{aligned}[t]&\phantom{={}}\color{black}678-(125\times3-278)\\&\color{answerred}=678-(375-278)\\&\color{answerred}=678-97\\&\color{answerred}=581\end{aligned}}
\MethodBlock{方法二：先处理括号或乘除，再去括号后计算}{加号后的括号去掉不变号，减号后的括号去掉要改变各项的加减号。}{\begin{aligned}[t]&\phantom{={}}678-(125\times 3-278)\\
&=678-(375-278)\\
&=678-375+278\\
&=303+278\\
&=581\end{aligned}}

\par\noindent

\QHeading{第3题　解决问题。}

\Problem{从A城到B城的路程为350千米。李叔叔坐小汽车从A城到B城，小汽车平均每小时行驶70千米；张老师坐客车从A城到B城，客车平均每小时行驶50千米。如果两人同时出发，张老师比李叔叔晚几小时到达？}
\par


\MethodTitle{方法一}
\Answer{\Formula{\begin{aligned}[t]&\phantom{={}}350\div50-350\div70\\&=7-5\\&=2\text{（小时）}\end{aligned}}\par 答：张老师比李叔叔晚2小时到达。}

\MethodTitle{方法二：看小汽车到达时客车还剩多少路}
\Answer{先求小汽车到达所用时间，再求同一时刻客车的剩余路程。}
\Answer{\textbf{小汽车所用时间}\par \Formula{\begin{aligned}[t]&\phantom{={}}350\div 70\\
&=5\text{（小时）}\end{aligned}}}
\Answer{\textbf{此时客车剩余路程}\par \Formula{\begin{aligned}[t]&\phantom{={}}350-50\times 5\\
&=350-250\\
&=100\text{（千米）}\end{aligned}}}
\Answer{\textbf{客车还要行驶的时间}\par \Formula{\begin{aligned}[t]&\phantom{={}}100\div 50\\
&=2\text{（小时）}\end{aligned}}}
\Answer{答：张老师比李叔叔晚2小时到达。}
